\documentclass{article}
\usepackage[utf8]{inputenc}
\usepackage[T1]{fontenc}
\usepackage[utf8]{inputenc}
\usepackage{lmodern}
\usepackage[a4paper, margin=1in]{geometry}
\usepackage{amsmath}
\usepackage{amssymb}
\usepackage{enumitem}
\usepackage{ragged2e}
\usepackage[spanish]{babel}
\setlist[itemize]{label=-}
\usepackage{booktabs}
\usepackage{xcolor}
\definecolor{primary}{RGB}{24, 43, 73}
\definecolor{secondary}{RGB}{71, 85, 105}
\definecolor{accentline}{RGB}{203, 213, 225}

\usepackage{hyperref}
\hypersetup{
	colorlinks=true,
	linkcolor=primary,
	urlcolor=primary,
	citecolor=primary
}
\large
\title{Tarea 1}
\begin{document}
	\begin{titlepage}
		\begin{center}
			\line(1,0){300}\\
			[0.65cm]
			\huge{\bfseries Tarea aplicada 1}\\
			\line(1,0){300}\\
			\textsc{\Large Álgebra lineal y Optimización}\\
			\today\\
			\textsc{\small Base de datos: Historial de precios de cierre ajustados (API Yahoo Finance) del S$\&$P 500 (2014-2024).}\\
			[5.5cm]     
		\end{center}
    % Datos de los 2 Integrantes en columnas paralelas
\begin{minipage}{0.95\textwidth}
	\begin{tabular}{@{}p{0.47\linewidth}@{\hspace{0.06\linewidth}}p{0.47\linewidth}@{}}
		\textbf{\color{primary}Integrante 1} & \textbf{\color{primary}Integrante 2} \\[0.3cm]
		\textbf{Nombre:} Jorge Ignacio Moreno Farías & \textbf{Nombre:} Miguel Ignacio Arias Meriño \\[0.15cm]
		\textbf{RUT:} 21.328.803-2 & \textbf{RUT:} [98.765.432-1] \\[0.15cm]
		\textbf{Correo:} \href{mailto:jorgmoreno@alumnos.uai.cl}{jorgmoreno@alumnos.uai.cl} & \textbf{Correo:} \href{mailto:estudiante2@correo.cl}{estudiante2@correo.cl} \\
		
	\end{tabular}
\end{minipage}

\vspace{1.8cm}



\end{titlepage}
		
	

	\section*{Introducción}
 
	
	En el presente trabajo se estudiará el clustering espectral como herramienta para analizar datos cuya estructura natural es una red. El clustering espectral tiene sus raíces en la teoría de grafos espectrales de la década de 1970 (con trabajos de Fiedler y Donath) y se popularizó en Machine Learning a principios de los 2000 gracias a investigadores como Shi, Malik, Ng, Jordan y Weiss. Entre sus principales ejemplos de uso destacan la segmentación de imágenes en visión por computadora, la detección de comunidades en redes biológicas y sociales, y el análisis de riesgo.\\
	
	A pesar de su eficacia, este método presenta limitaciones, como su alta sensibilidad a valores atípicos (outliers) que distorsionan el Laplaciano, y el alto costo computacional al calcular vectores propios en matrices de gran escala. La motivación principal de este estudio en Ciencia de Datos, y en particular en el análisis de riesgo de mercados financieros, es poder identificar agrupaciones o comunidades basándonos puramente en la fuerza de sus conexiones, es decir, las correlaciones, demostrando por qué una representación matricial y espectral de una red es útil para filtrar el ruido y capturar la topología del riesgo sistémico.
	
	\section*{Desarrollo}
	\subsection*{Diseño experimental}
	Se evaluaron sistemáticamente dos decisiones metodológicas centrales:
	
	1. \textbf{Laplaciano No Normalizado vs. Normalizado Simétrico:} 
	Al utilizar el Laplaciano estándar \(L = D - A\), el algoritmo minimizó el corte absoluto \(\text{cut}(A,B)\). Debido al umbral de similitud (\(\rho \ge 0,40\)), la red presentó dos activos desconectados (\texttt{CHRD} y \texttt{MUSA}). Como se observa en la Figura~\ref{fig:grafo_fallido}, el modelo cayó en una partición trivial: aisló estos nodos individuales y agrupó al 98\% del mercado en un único clúster gigante, arrojando un Adjusted Rand Index (ARI) nulo de $0,021$. Al implementar el Laplaciano Normalizado Simétrico \(L_{\text{sym}} = I - D^{-1/2}AD^{-1/2}\) junto a la normalización de filas a norma unitaria (algoritmo NJW), se forzó un Normalized Cut penalizado por volumen, logrando la separación balanceada de los sectores.
	
	2. \textbf{Validación frente a Scikit-Learn y Sensibilidad de Umbral:}
	Se contrastó la implementación manual matricial frente a la función de alto nivel \texttt{SpectralClustering} de Scikit-Learn, obteniendo asignaciones de clústeres idénticas (\(\text{ARI} = 0,835\)), lo que verifica la reproducibilidad numérica del pipeline. Asimismo, se analizó la sensibilidad del umbral: umbrales menores a $0,30$ saturan la red de ruido estocástico difuminando las fronteras sectoriales, mientras que umbrales mayores a $0,55$ fracturan el grafo en múltiples componentes desconectadas.
	
	
	\section*{Desarrollo: Parte I. Fundamentos matemáticos}
	
	\subsection*{Pregunta 1. Una red como objeto matricial}
	\begin{enumerate}[label=\alph*)]
		\item \textbf{Información que contienen \(A\), \(D\) y \(L\):} La matriz de adyacencia \(A\) nos dice exactamente qué nodos están conectados y con qué fuerza (pesos). La matriz diagonal de grados \(D\) resume el peso total de las conexiones de cada nodo, es decir, qué tan conectado o central es ese nodo en toda la red. La matriz Laplaciana \(L = D - A\) combina ambas para capturar la estructura completa del grafo, indicando cómo fluye la información y dónde están las tensiones o fronteras naturales de la red.
		
		\item \textbf{Interpretación del vector \(x \in \mathbb{R}^n\):} Si la red es un mercado, el vector \(x\) es una función que le asigna un número (como una ``etiqueta'' o ``estado'') a cada acción de la red. Cada componente \(x_i\) es el valor que toma el nodo \(i\).
		
		\item \textbf{Espacio vectorial de funciones:} El conjunto de estas funciones forma un espacio vectorial porque si tenemos dos asignaciones (vectores) distintas \(x, y \in \mathbb{R}^n\), podemos sumarlas nodo a nodo (\(z_i = x_i + y_i\)) y el resultado sigue siendo una asignación válida para los nodos de la red. De igual forma, podemos ponderar una función por un escalar (ej. $2x$), cumpliendo así las propiedades de combinación lineal.
		
		\item \textbf{Demostración de que \(L\) es simétrica:} Como el grafo es no dirigido, las conexiones son bidireccionales con el mismo peso, por lo que \(A = A^\top\). La matriz \(D\) es diagonal, por lo que trivialmente es simétrica (\(D = D^\top\)). Al ser \(L = D - A\), aplicando la transpuesta tenemos:
		\[
		L^\top = (D - A)^\top = D^\top - A^\top = D - A = L.
		\]
		Por lo tanto, \(L\) es simétrica.
	\end{enumerate}
	
	\subsection*{Pregunta 2. La forma cuadrática del Laplaciano}
	\textbf{Demostración de la identidad:} Desarrollamos la expresión matricial \(x^\top L x\):
	\[
	x^\top L x = x^\top (D - A) x = x^\top D x - x^\top A x
	\]
	Sabiendo que \(D\) es diagonal, \(x^\top D x = \sum_i d_i x_i^2\). Y para \(A\), \(x^\top A x = \sum_{i,j} a_{ij} x_i x_j\). Reescribiendo \(\sum_i d_i x_i^2\) y sabiendo que \(d_i = \sum_j a_{ij}\):
	\[
	x^\top D x = \sum_i \left( \sum_j a_{ij} \right) x_i^2 = \frac{1}{2} \sum_{i,j} a_{ij} x_i^2 + \frac{1}{2} \sum_{i,j} a_{ij} x_j^2
	\]
	Restando el término de \(A\):
	\[
	x^\top L x = \frac{1}{2} \sum_{i,j} a_{ij} x_i^2 - \sum_{i,j} a_{ij} x_i x_j + \frac{1}{2} \sum_{i,j} a_{ij} x_j^2
	\]
	Agrupando términos y reconociendo el cuadrado de binomio:
	\[
	x^\top L x = \frac{1}{2} \sum_{i,j} a_{ij} (x_i^2 - 2x_i x_j + x_j^2) = \frac{1}{2} \sum_{i,j} a_{ij} (x_i - x_j)^2
	\]
	
	\begin{enumerate}[label=\alph*)]
		\item \textbf{\(L\) es semidefinida positiva:} Como los pesos \(a_{ij}\) son no negativos (\(a_{ij} \ge 0\)) y un cuadrado \((x_i - x_j)^2\) siempre es mayor o igual a cero, la suma total siempre es \(\ge 0\) para cualquier vector \(x\).
		
		\item \textbf{Interpretación geométrica:} Representa el ``costo'' o ``tensión'' total de la asignación sobre la red. Penaliza las diferencias de valor entre nodos conectados.
		
		\item \textbf{Nodos fuertemente conectados con valores similares:} Si dos nodos tienen una conexión fuerte (\(a_{ij}\) grande) y se les asigna un valor similar (\(x_i \approx x_j\)), la diferencia al cuadrado es casi cero, por lo que no aportan penalización (tensión) a la suma total.
		
		\item \textbf{\(L\mathbf{1} = 0\) y su interpretación:} El vector \(\mathbf{1}\) es un vector de puros unos. \(L\mathbf{1} = D\mathbf{1} - A\mathbf{1}\). \(D\mathbf{1}\) es un vector donde la fila \(i\) tiene el grado \(d_i\). \(A\mathbf{1}\) es un vector donde la fila \(i\) es la suma de la fila \(i\) de \(A\), que por definición es \(d_i\). Por ende, \(D\mathbf{1} - A\mathbf{1} = 0\). Geométricamente, significa que si asignamos el mismo valor a todos los nodos (un solo grupo general), el ``costo de corte'' es cero.
	\end{enumerate}
	
	\subsection*{Pregunta 3. Matriz de incidencia y descomposición SVD}
	\begin{enumerate}[label=\alph*)]
		\item \textbf{Demostración \(L = B^\top B\) para red no ponderada:} La matriz \(B\) tiene filas por cada arista \(e = (u, v)\) y columnas por nodo. Para una arista \(e\), \(B_{eu} = 1\) y \(B_{ev} = -1\), el resto $0$. Al hacer \((B^\top B)_{ij}\) multiplicamos la columna \(i\) por la columna \(j\) de \(B\). Si \(i = j\) (la diagonal), multiplicamos la columna de un nodo por sí misma, sumando $1$ por cada arista que incide en él, lo que da el grado \(d_i\). Si \(i \ne j\), el producto es \(-1\) si existe una arista entre ellos, y $0$ si no. Esto es exactamente \(L\).
		
		\item \textbf{Modificación para incorporar pesos:} Para una red ponderada, se multiplica cada fila de la matriz de incidencia \(B\) (correspondiente a una arista \(e\)) por la raíz cuadrada del peso de esa arista, es decir, \(\sqrt{w_e}\).
		
		\item \textbf{Descomposición SVD:} Sabiendo que \(B = U \Sigma V^\top\):
		\[
		L = B^\top B = (U \Sigma V^\top)^\top (U \Sigma V^\top) = V \Sigma^\top U^\top U \Sigma V^\top
		\]
		Como \(U\) es ortogonal, \(U^\top U = I\). Por lo tanto, \(L = V \Sigma^\top \Sigma V^\top\).
		
		\item \textbf{Relación entre SVD y valores propios:} De la ecuación anterior, \(L = V \Sigma^2 V^\top\). Esto tiene exactamente la forma de la diagonalización de \(L\). Así, los vectores singulares derechos de \(B\) (las columnas de \(V\)) son exactamente los vectores propios de \(L\), y los valores propios de \(L\) son los cuadrados de los valores singulares de \(B\).
	\end{enumerate}
	
	\subsection*{Pregunta 4. El espectro y la conectividad de la red}
	\begin{enumerate}[label=\alph*)]
		\item \textbf{Uso del teorema espectral:} El teorema espectral establece que toda matriz real y simétrica es diagonalizable ortogonalmente. Como demostramos en (1d) que \(L\) es simétrica, existe una base ortonormal de \(\mathbb{R}^n\) compuesta por sus vectores propios.
		
		\item \textbf{Componentes conexas y núcleo de \(L\):} Si hay \(k\) componentes conexas, no hay aristas entre ellas. Podemos ordenar \(L\) en forma de bloques. Cada bloque es un Laplaciano válido que tiene un autovalor $0$ con el vector de unos para esa componente. Así, tenemos al menos \(k\) vectores propios de valor $0$ que son linealmente independientes (el núcleo tiene dimensión \(k\)).
		
		\item \textbf{Multiplicidad de \(\lambda = 0\):} La multiplicidad algebraica del autovalor $0$ es exactamente igual al número de componentes conexas totalmente aisladas en el grafo.
		
		\item \textbf{Información de \(\lambda_2\):}
		\begin{itemize}
			\item \textbf{Igual a 0:} Hay al menos 2 componentes desconectadas (la red está partida en dos).
			\item \textbf{Positivo pero pequeño:} La red es conexa, pero existe un ``cuello de botella'' muy débil que casi la divide en dos.
			\item \textbf{Relativamente grande:} La red está densamente conectada, sin fronteras claras fáciles de cortar.
		\end{itemize}
		
		\item \textbf{Vector de Fiedler:} Sus coordenadas asignan valores positivos a un grupo de nodos y negativos a otro. Al minimizar la tensión, el vector ubica las diferencias mayores precisamente sobre los enlaces más débiles de la red, revelando el corte natural en dos comunidades.
	\end{enumerate}
	
	\subsection*{Pregunta 5. Partición de una red como problema de optimización}
	\begin{enumerate}[label=\alph*)]
		\item \textbf{Representación matemática:} Un vector \(x \in \mathbb{R}^n\) donde \(x_i = 1\) si el nodo pertenece al grupo \(A\), y \(x_i = -1\) si pertenece al grupo \(B\).
		
		\item \textbf{Función del peso total:} El peso del corte (\(\text{Cut}\)) es la suma de los pesos \(a_{ij}\) de todas las aristas que conectan nodos con diferente signo en \(x\).
		
		\item \textbf{Cortes desbalanceados:} Si solo minimizamos el peso del corte, el algoritmo suele separar un solo nodo aislado del resto (porque cortar una sola conexión es muy barato), lo que no es útil en la práctica.
		
		\item \textbf{Ratio Cut / Normalized Cut:} Para evitar cortes triviales, se introduce un denominador que penaliza grupos muy pequeños. Ratio Cut divide el costo por la cantidad de nodos de cada grupo, y Normalized Cut lo divide por la suma de los grados (volumen) de cada grupo.
		
		\item \textbf{Relajación:} El problema discreto de usar valores exactos \(\{1, -1\}\) es NP-hard. Relajamos la restricción permitiendo que \(x\) tome cualquier valor real continuo, lo que transforma el problema discreto en minimizar la forma cuadrática continua \(x^\top L x\).
		
		\item \textbf{Solución del problema continuo:} Por el Teorema de Rayleigh-Ritz, minimizar \(x^\top L x\) sujeto a \(\|x\| = 1\) es buscar los vectores propios. La restricción \(x^\top \mathbf{1} = 0\) obliga a que la solución sea ortogonal al primer vector propio (asociado a \(\lambda_1 = 0\)). Por ende, el mínimo siguiente es el segundo valor propio \(\lambda_2\), y la solución es su respectivo vector propio.
		
		\item \textbf{Conversión a partición de nodos:} Para volver al mundo discreto, podemos mirar el signo de cada coordenada del vector propio. Si \(x_i > 0\) va al grupo \(A\), y si \(x_i \le 0\) va al grupo \(B\). Para problemas más complejos, se aplica \(k\)-means a esos valores.
	\end{enumerate}
	
	\subsection*{Pregunta 6. De la bipartición al clustering espectral}
	\begin{enumerate}[label=\alph*)]
		\item \textbf{Fila \(i\) de \(U\):} Representa el nuevo ``perfil espectral'' (embedding) del nodo \(i\), extrayendo sus características estructurales en la red.
		
		\item \textbf{Interpretación en \(\mathbb{R}^k\):} Hemos transformado el nodo de ser simplemente un vértice del grafo a ser una coordenada espacial en un espacio de \(k\) dimensiones geométricas.
		
		\item \textbf{Distancia en el agrupamiento:} En este nuevo espacio, si dos puntos (nodos) están a poca distancia euclidiana, significa que comparten el mismo patrón de conexiones en el grafo original (pertenecen a la misma comunidad estructural).
		
		\item \textbf{Uso de \(k\)-means:} Como ahora los nodos son simples puntos en \(\mathbb{R}^k\), podemos aplicar el clásico algoritmo de \(k\)-means iterativo sobre las filas de la matriz \(U\) para agrupar los puntos más cercanos geométricamente en \(k\) clusters distintos.
		
		\item \textbf{Comparación conceptual:} Usar características originales aísla cada nodo sin ver el contexto; usar la matriz de adyacencia directa es ruidoso y de alta dimensión; usar la representación espectral filtra el ruido y reduce la dimensionalidad para destacar la macrotopología de las comunidades.
		
		\item \textbf{Elección usando eigengap:} El eigengap es un salto grande en el valor entre autovalores consecutivos (\(\lambda_{k+1} - \lambda_k\)). Buscar este salto nos sugiere el número \(k\) natural de grupos densamente conectados en la red.
		
		\item \textbf{Dificultades del clustering espectral:}
		\begin{enumerate}[label=\arabic*.]
			\item Muy sensible a valores atípicos y ruidosos, donde una arista con peso excesivamente anómalo distorsiona el Laplaciano.
			\item Si los grupos tienen densidades o tamaños muy diferentes (incluso usando Laplaciano normalizado), el \(k\)-means posterior puede fallar al definir los centroides correctamente.
		\end{enumerate}
	\end{enumerate}
	
	\section*{Parte II. Aplicación a una red real}
	
	\subsection*{5.1. Problema y datos}
	La base de datos utilizada corresponde al historial de precios de cierre ajustados de un universo de aproximadamente 145 acciones pertenecientes al S\&P 500, abarcando el periodo 2014--2024. Los datos fueron obtenidos mediante la API pública de Yahoo Finance.
	
	En esta red, los nodos representan acciones individuales de cuatro sectores económicos reales (Tecnología, Finanzas, Salud y Energía). Las aristas tienen pesos y la red es no dirigida; el peso de la conexión entre dos nodos está dado por la matriz de correlación de Pearson de sus retornos diarios.
	
	La pregunta de análisis que guía esta aplicación es: \textit{¿Logra el clustering espectral descubrir comunidades de riesgo estructural en el mercado que coincidan con las clasificaciones sectoriales tradicionales (GICS), basándose puramente en las fluctuaciones de precio y sin conocer las etiquetas previas?}
	
	El clustering espectral es pertinente para este problema porque las correlaciones financieras forman un sistema complejo y denso (con alto ruido diario) donde los algoritmos tradicionales de agrupamiento fallan por la maldición de la dimensionalidad. El espectro de la matriz Laplaciana permite filtrar este ruido y capturar la topología global del riesgo sistémico.
	
	\subsection*{5.2. Construcción y preparación de la red}
	Para construir la red, se tomaron las siguientes decisiones metodológicas:
	\begin{itemize}
		\item \textbf{Filtrado temporal y de nodos faltantes:} Se utilizaron 10 años de historia bursátil para garantizar robustez estadística. Se eliminaron automáticamente los activos que no contaban con el historial completo para evitar distorsiones matemáticas en la correlación.
		\item \textbf{Tratamiento de pesos y umbral:} Como la correlación asume valores en \([-1, 1]\) y la matriz de adyacencia requiere pesos no negativos que reflejen similitud, se aplicó un umbral estricto. Toda correlación menor a $0{,}40$ (y correlaciones negativas) fue forzada a $0$. Esto limpia la red de enlaces débiles (ruido).
		\item \textbf{Eliminación de loops:} Se forzó la diagonal de la matriz de adyacencia a $0$, dado que la autocorrelación de un activo ($1{,}0$) genera bucles redundantes que no aportan a la topología de la partición.
	\end{itemize}
	
	\subsection*{5.3. Implementación del método espectral}
	El proceso fue implementado en Python utilizando bibliotecas numéricas (\texttt{numpy}, \texttt{scipy.linalg}, \texttt{scikit-learn}). A partir de la matriz de adyacencia (\(A\)), se calculó la matriz de grados (\(D\)) y el Laplaciano estándar (\(L = D - A\)). Posteriormente, se calcularon los valores y vectores propios de \(L\). Utilizando el gráfico del espectro para observar el eigengap, y guiados por el conocimiento previo de los \(k = 4\) sectores macro introducidos en la muestra, se construyó la matriz de representación espectral \(U \in \mathbb{R}^{n \times 4}\). Finalmente, se aplicó el algoritmo \(k\)-means sobre las filas normalizadas de esta matriz.
	
	\subsection*{5.4. Diseño experimental}
	El experimento central consistió en evaluar la sensibilidad del algoritmo frente a la elección del tipo de Laplaciano (Laplaciano estándar no normalizado vs. Laplaciano Normalizado Simétrico). Al utilizar el Laplaciano estándar (\(L = D - A\)), el algoritmo intentó minimizar estrictamente el peso del corte. Debido al umbral de correlación impuesto ($0{,}40$), la red presentaba nodos pobremente conectados en su periferia. El algoritmo optó por aislar nodos individuales (creando componentes conexas triviales) y agrupó al \(\sim 98\%\) del mercado en un único clúster gigante, fallando en su tarea de discriminar sectores. Para solucionar esto, se aplicó el Laplaciano Normalizado (\(L_{\text{norm}} = I - D^{-1/2} A D^{-1/2}\)). Esto obligó al modelo a realizar un Normalized Cut, penalizando fuertemente los cortes desbalanceados e induciendo la búsqueda de grupos con volumen equiparable.
	
	\subsection*{5.5. Resultados e interpretación}
	Los resultados con el Laplaciano Normalizado fueron sumamente exitosos. La visualización topológica de la red demostró que el clustering espectral logró reconstruir ``a ciegas'' los sectores macroeconómicos de la bolsa:
	\begin{itemize}
		\item El algoritmo identificó y separó limpiamente a las empresas energéticas (ej. XOM, CVX) debido a su alta covarianza interna.
		\item Separó al bloque financiero y al tecnológico en comunidades densas distintas.
		\item Se observaron ``nodos fronterizos'' (ej. FSLR, ENPH) que qued3aron periféricos a los clústeres principales. Aunque legalmente pertenecen a sectores como Energía o Tecnología, su dinámica de retornos mostró un comportamiento atípico (energías renovables), demostrando que la red estructural captura el riesgo real latente más allá de la etiqueta formal.
	\end{itemize}
	
	\subsection*{Resultados e interpretación}
	La búsqueda del Eigengap en el espectro del Laplaciano (Figura~\ref{fig:eigengap}) reveló tres autovalores iniciales nulos (\(\lambda_1 = \lambda_2 = \lambda_3 = 0\)), confirmando teóricamente la presencia de 3 componentes conexas en el grafo. Tras \(\lambda_4\), se observa un salto marcado hacia \(\lambda_5\), justificando la elección de \(k=4\).
	
	Al proyectar las comunidades en el grafo topológico (Figura~\ref{fig:grafo_normalizado}), el algoritmo reconstruyó con alta precisión los cuatro sectores reales:
	\begin{itemize}
		\item \textbf{Energía:} Formó un bloque denso e independiente (\texttt{XOM}, \texttt{CVX}, \texttt{SLB}), producto de su alta covarianza interna ligada al precio del crudo.
		\item \textbf{Finanzas:} Agrupó a la banca y aseguradoras (\texttt{JPM}, \texttt{GS}, \texttt{BAC}) con mínima dispersión.
		\item \textbf{Tecnología y Salud:} Quedaron desacoplados en macro-comunidades estables.
		\item \textbf{Nodos Fronterizos:} Activos como First Solar (\texttt{FSLR}) y Enphase Energy (\texttt{ENPH}) quedaron en la frontera entre Energía y Tecnología. Aunque pertenecen formalmente al sector energético, su dinámica de retornos responde al ciclo de innovación, semiconductores y tasas de interés, demostrando que la representación espectral captura el riesgo sistémico subyacente más allá de la clasificación burocrática.
	\end{itemize}
	
	\begin{table}[h!]
		\centering
		\small
		\begin{tabular}{lcccc}
			\hline
			\textbf{Modelo / Laplaciano} & \textbf{ARI} & \textbf{NMI} & \textbf{Silueta} & \textbf{Modularidad (\(Q\))} \\ \hline
			Laplaciano Estándar (\(L\)) & 0,021 & 0,045 & -0,120 & 0,031 \\
			Laplaciano Normalizado (\(L_{\text{sym}}\)) & \textbf{0,835} & \textbf{0,791} & \textbf{0,488} & \textbf{0,468} \\
			Benchmark Scikit-Learn & \textbf{0,835} & \textbf{0,791} & \textbf{0,488} & \textbf{0,468} \\ \hline
		\end{tabular}
		\caption{Métricas cuantitativas de evaluación de la partición.}
		\label{tab:metricas}
	\end{table}
	\section*{Aplicación}
	%%%%%%%%%%%%%%%%%%%%%%%
	
	%%%%%%%%%%%%%%%%%%%%%%
	\section*{Conclusión}
	\section{Conclusiones}
	En este trabajo se demostró la efectividad del clustering espectral como herramienta para la detección no supervisada de comunidades de riesgo en mercados financieros complejos. A través del análisis del espectro Laplaciano, fue posible sortear la maldición de la dimensionalidad inherente a las series de tiempo bursátiles, transformando un problema combinatorio NP-hard en una relajación continua elegante basada en autovectores.\\
	
	Empíricamente, se constató que la formulación no normalizada del Laplaciano es vulnerable a particiones triviales cuando existen activos con baja conectividad periférica, haciendo indispensable el uso del Laplaciano Normalizado Simétrico (\(L_{\text{sym}}\)) y la proyección a la esfera unitaria propuesta por Ng-Jordan-Weiss. La red resultante no solo reconstruyó fielmente los cuatro macro-sectores del S\&P 500 sin supervisión previa, sino que identificó nodos fronterizos cuya covarianza real desafía las clasificaciones administrativas tradicionales, aportando información relevante para la diversificación y el control de riesgo en portafolios de inversión.
	
	\subsection*{5.6. Limitaciones y extensiones}
	Una limitación detectada es la sensibilidad extrema al umbral de correlación; un umbral muy alto desconecta el grafo (generando ceros en los autovalores), mientras que uno muy bajo lo vuelve excesivamente denso, difuminando el eigengap. Asimismo, la correlación de Pearson asume relaciones lineales estáticas, ignorando los cambios de régimen de mercado durante crisis financieras.
	
	Como trabajo futuro, se propone aplicar clustering espectral sobre redes temporales (ventanas móviles de correlación) para observar si las comunidades evolucionan dinámicamente antes o después de una crisis, o utilizar Graph Neural Networks (GNNs) para integrar tanto la topología de la red como atributos fundamentales de cada empresa (balances, capitalización) en el proceso de agrupación.
	
	
	
	\section*{Referencias}

	\begin{enumerate}
		\item von Luxburg, U. (2007). \textit{A tutorial on spectral clustering}. Statistics and Computing, 17(4), 395-416.
		\item Shi, J., \& Malik, J. (2000). \textit{Normalized cuts and image segmentation}. IEEE Transactions on Pattern Analysis and Machine Intelligence, 22(8), 888-905.
		\item Ng, A. Y., Jordan, M. I., \& Weiss, Y. (2001). \textit{On spectral clustering: Analysis and an algorithm}. Advances in Neural Information Processing Systems (NIPS), 14, 849-856.
		\item Fiedler, M. (1973). \textit{Algebraic connectivity of graphs}. Czechoslovak Mathematical Journal, 23(2), 298-305.
		\item Newman, M. E. (2006). \textit{Modularity and community structure in networks}. Proceedings of the National Academy of Sciences, 103(23), 8577-8582.
		\item Mantegna, R. N. (1999). \textit{Hierarchical structure in financial markets}. The European Physical Journal B, 11(1), 193-197.
		\item Hagberg, A., Swart, P., \& S Chult, D. (2008). \textit{Exploring network structure, dynamics, and function using NetworkX}. Los Alamos National Laboratory (LANL).
		\item Pedregosa, F., et al. (2011). \textit{Scikit-learn: Machine learning in Python}. Journal of Machine Learning Research, 12, 2825-2830.
	\end{enumerate}
	\section*{Código}
	Aquí podrá encontrar el código que utilizamos	\href{https://github.com/troncos22/Tareas_2026_2}{github}
	\section*{Uso de IA}
	El uso de Ia en su mayoría fue para revisar documentación de las librerías utilizadas y de documentación de LaTex. Verificación de rubricas y de si faltaba algo tanto en el código, informe y presentación. También para la resolución de problemas de compilación de código Python y LaTex. 
	\subsection*{Promts}
		\begin{itemize}
			\item $"$ween día oye necesito ayuda
			con mi compañero desarrollamos este código y quería saber que le pude faltar a partir de esta lista que hice anteriormente contigo
			https://share.gemini.google/kKcNXgEE4FqM $"$
			\item $"$esta es la tarea, a partir del código anterior qué le falta?$"$
			\item $"$y a éste documento qué le falta?$"$
			\item $"$ quiero que el texto aparezca justificado $"$
		\end{itemize}
	\subsection*{main.cpp}

	
\end{document}











